% year: תשע"ז
% type: מבחן
% semester: ב
% moed: ב
% source: exams/מבחנים/תשעז/מועד ב תשעז סמסטר ב.pdf
% number: 8
% points: 10
% categories: definite-integrals, true-false

\begin{question}
הוכיחו או הפריכו: אם הפונקציה $|f(x)|$ אינטגרבילית לפי רימן בקטע $[a,b]$ אז גם הפונקציה $f(x)$ אינטגרבילית בקטע זה.
\end{question}

\begin{hint}
חפשו פונקציה שמקבלת רק את הערכים $1$ ו-$-1$ באופן "צפוף", בדומה לפונקציית דירכלה.
\end{hint}

\begin{solution}
\textbf{הטענה אינה נכונה.} (הכיוון ההפוך נכון: אם $f$ אינטגרבילית אז גם $|f|$.)

דוגמה נגדית (עבור $a<b$):
\[ f(x)=\begin{cases}1 & x\in\Q\\ -1 & x\notin\Q\end{cases}\qquad x\in[a,b]. \]
אז $|f(x)|=1$ לכל $x$, פונקציה קבועה, ולכן אינטגרבילית ו-$\int_a^b|f|=b-a$.

לעומת זאת $f$ אינה אינטגרבילית לפי רימן: תהי $P$ חלוקה כלשהי של $[a,b]$. בכל תת-קטע $[x_{i-1},x_i]$ (באורך חיובי) יש גם מספר רציונלי וגם מספר אי-רציונלי (צפיפות $\Q$ ו-$\R\setminus\Q$ ב-$\R$), ולכן $M_i=\sup f=1$ ו-$m_i=\inf f=-1$. מכאן
\[ U(f,P)=\sum M_i\Delta x_i=b-a,\qquad L(f,P)=\sum m_i\Delta x_i=-(b-a), \]
כך ש-$U(f,P)-L(f,P)=2(b-a)$ לכל חלוקה, ואינו יכול להיות קטן מ-$\eps$ כרצוננו. לפי קריטריון רימן (או: האינטגרל העליון $b-a$ שונה מהתחתון $-(b-a)$), $f$ \textbf{אינה} אינטגרבילית.
\end{solution}
